%%%%%%%%%%%%%%%%%%%%%%%%%%%%%%%%%%%%%%%%%%%%%%%%%
%================================================
% Sec. 7 Discussion and Sec. 8 Conclusions (draft 2026-10-28, replaces the WP 0.1 skeleton).
% Outline: numerical_tests_plan_2026-10-22.md Sec. 4; agent_handoff_2026-10-26.md Sec. 4.
% All numbers are those of Sec. 6 (numerical_results.tex), except the L3-0 sizing figures in
% Sec. 7.3 (marked [L3-0]), which come from docs/L3_implementation_report.md.
% L3-3, L4-4 and L4-8 are pending: their sentences are placeholders, marked [PENDING].
% Text in red until the final pass.
%================================================
{\color{red}
\section{Discussion}\label{sec:discussion}

\subsection{Validity of the local theory}\label{subsec:disc_validity}

The local effective-medium model of Section~\ref{sec:dispersion_relation} rests on three parameters. The scale separation $\varepsilon_s=d_s/\lambda$, equation~\eqref{eq:homogenization_scale_separation}, measures the spacing against the wavelength. The lattice-compliance number $\Pi_L$, equation~\eqref{eq:PiL}, measures the near-field compliance of the structure over one cell against the resonator spring. The fluid-diluted mass ratio $\mu(k_c)$, equation~\eqref{eq:lat_mu}, measures the resonator mass against the host inertia, fluid included, at the crossing wavenumber. The finite element results of Section~\ref{sec:numerical_results} show what each parameter controls.

\paragraph{Scale separation.} At fixed loading, the hybridisation of the discrete line lattice approaches the local value $\sqrt{\mu(k_c)}$ with a correction of order $\varepsilon_s^{4.5}$: $2.7\times10^{-4}$ at $\varepsilon_s=\tfrac14$ and $6.2\times10^{-7}$ at $\varepsilon_s=\tfrac1{16}$ (Section~\ref{subsubsec:num_L1_8}). The local dispersion relation is therefore accurate well before $\varepsilon_s\ll1$. The exception is a lattice that is coarse compared with the evanescent decay length, e.g.\ $\omega_r=2~\si{\radian\per\second}$, where the correction is $4.9\times10^{-5}$ at $\varepsilon_s=\tfrac1{16}$. Small $\varepsilon_s$ does not, however, make the local model complete. The discrete lattice has a complete stop band, whose two edges are set by the attachment pattern (Section~\ref{subsubsec:num_L1_5}). At fixed loading its width vanishes as $a^4$, with a finite element slope of $4.01$ (Section~\ref{subsubsec:num_L1_6}). The homogenised model has no complete stop band at all: its finite element discretisation transmits at every sampled frequency except $\omega_r$ (Section~\ref{subsubsec:num_L1_7}). Truncating the local branches at the zone boundary does not recover the gap. It misplaces both edges and underestimates the width by a factor of $1.8$--$1.9$. For the parameters of Table~\ref{tab:dispersion_test_parameters} at $a=1~\si{\meter}$, $\varepsilon_s=0.44$ and $0.31$ at the crossing. This is the regime where local resonance and Bragg scattering interact, and where only the discrete model applies.

\paragraph{Near-field loading.} The zone-centre edge of the discrete lattice is the resonance renormalised by the lattice near field, $\omega_\Gamma\simeq\omega_r/\sqrt{1+\Pi_L}$, equation~\eqref{eq:lat_gamma}. The estimate holds for $2\pi/a>k_c$, where $\Pi_L>0$. There its error is below $0.44\,\Pi_L^2$ and decreases with a slope of $3.13$--$3.16$ in $\Pi_L$, because the coefficient of $\Pi_L^2$ is itself proportional to $\Pi_L$ when bending dominates (Section~\ref{subsubsec:num_L1_7}). For $2\pi/a<k_c$, $\Pi_L<0$ and the estimate fails. Since $\Pi_L$ does not depend on the wavelength, a small $\varepsilon_s$ is necessary but not sufficient for the local model: stiff resonators on a flexible host shift the resonance even on a deeply subwavelength lattice.

\paragraph{Fluid dilution.} The hybridisation of the floating lattice is $\sqrt{\mu(k_c)}$, between $0.064$ and $0.20$ in the cases of Section~\ref{subsubsec:num_L1_8}, and not the in-vacuo mass ratio, which is one. Depth enters only through $k_c$: the splittings at $H=2$, $10$ and $100~\si{\meter}$ agree to $1.3\times10^{-5}$. On a finite structure the cut-off of the added mass at $k\sim\pi/L$ reduces the hybridisation as $L^{-1/2}$ in deep water and towards $L^{-1}$ in shallow water (Section~\ref{subsubsec:num_L2_5}). These exponents are confirmed, but the continuum estimate $\sqrt{\mu_B}$ underestimates the splitting by a factor of $1.7$--$1.8$ in deep and intermediate water. The wet modal mass of the bare mode gives the splitting to $1$--$6.5\%$. The same prefactor, of order $1.5$, applies to the estimate~\eqref{eq:mode_dilution} in the vertical section; its three-dimensional check is being completed [PENDING: L4-8].

\paragraph{Graded arrays.} For graded arrays the local model locates the absorption but does not give its spatial width (Section~\ref{subsec:num_graded}). For $\zeta\le0.02$ the peak of the dissipation lies within $\Delta x/2$ of the turning point $x^*$, equations~\eqref{eq:x_star_linear}--\eqref{eq:x_star_exp}, in all cases but one marginal case ($0.52\,\Delta x$). At $\zeta=0.05$ the wave is absorbed before it reaches $x^*$, and the peak moves upstream by $0.6$--$1.0\,\Delta x$. The width is $c\,2\zeta\,\omega|dx^*/d\omega|$, with $c$ between about $1$ and $3$, rising with the opacity of the array. The term $\sqrt{\mu}$ in equation~\eqref{eq:footprint_general} bounds the frequency window of the dispersion branch, not the spatial width. These tolerances are applied at $a=0.25~\si{\meter}$. At $a=1~\si{\meter}$, $\varepsilon_s$ reaches $0.36$ and the local law is not expected to hold.

In summary, the local model gives the dispersion branches and the hybridisation strength when $\varepsilon_s$ is at most about $\tfrac14$ (correction $2.7\times10^{-4}$ there) and $\Pi_L\ll1$, and it locates the absorption along a graded array. It does not give the complete stop band, the near-field shift of the resonance, the spatial width of graded absorption or, by itself, the reflection and transmission of finite arrays. Sections~\ref{sec:discrete_lattice} and~\ref{sec:numerical_results} supply these.

\subsection{Design consequences}\label{subsec:disc_design}

\paragraph{Resonator mass.} The resonator mass per unit area is to be compared with the mass of a fluid layer of thickness $1/k_c$, equation~\eqref{eq:dilution_design_rule}, not with the plate mass. The discrete lattice confirms this rule (Section~\ref{subsubsec:num_L1_8}). On a finite plate the relevant inertia is the wet modal mass of the target mode, which is smaller than the plane-wave added mass at the modal wavenumber because the fluid can flow round the free edges (Section~\ref{subsubsec:num_L2_5}). The hybridisation of a uniform population then decreases with the length of the plate. On a short plate the positions of the resonators also matter: with two resonators on a plate of $2.3~\si{\meter}$, both lie close to the nodes of the free--free mode and the splitting is small.

\paragraph{Damping.} Critical coupling, $\gamma_d=\gamma_r$, gives the absorbed fraction $\tfrac12$ at every frequency tested (Section~\ref{subsubsec:num_L1_3b}). The damping ratio that achieves it collapses at low frequency, from $9.26\times10^{-3}$ at $\omega_r=10~\si{\radian\per\second}$ to $4.35\times10^{-5}$ at $\omega_r=1~\si{\radian\per\second}$ ($T_\rho=0$), and from $6.37\times10^{-3}$ to $4.14\times10^{-5}$ ($T_\rho=9.8~\si{\meter\cubed\per\second\squared}$). At $\omega_r=1~\si{\radian\per\second}$ the radiative quality factor is $Q_\mathrm{rad}=1.15\times10^4$. A single resonator of fixed mass therefore becomes an extremely narrow-band absorber in the ocean-wave range, and the fluid, not the structure, sets this limit. The same dilution sets the loading needed for a target hybridisation, which grows as $\omega_r^{-2}$ at low frequency, equation~\eqref{eq:mass_grading_rule}.

\paragraph{Asymmetry.} A single resonator on a symmetric host absorbs at most one half of the incident power. On a finite structure the bound~\eqref{eq:lat_absorbed_B} is set by the ratio of the waves that the attachment radiates to the two sides (Section~\ref{subsubsec:num_L2_4}). Along a free plate it ranges from $0.088$ to $0.912$ ($L_p=5~\si{\meter}$) and from $0.10$ to $0.90$ ($L_p=20~\si{\meter}$). At $x_1=L_p/8$ on the short plate a single critically damped resonator absorbs $91\%$ of the incident power. The bound does not grow monotonically towards the edge; it follows the phase of the wet mode at the attachment. The best position is found from $G_B$ alone, without any resonator computation. The critical damping ratio, $0.003$ to $0.013$, is proportional to $\gamma_r(x_1)$, which changes by a factor of up to six between the nodes and the antinodes of the wet modes (Section~\ref{subsubsec:num_L2_1}).

\paragraph{Long and graded arrays.} The attenuation per unit length of a long floating array is bounded by approximately $\kappa_1\simeq\pi/(2H)$, the decay rate of the slowest evanescent fluid mode, whatever the resonators (Section~\ref{subsubsec:num_L1_4}). Adding resonators beyond this limit does not make the array more opaque. Along a graded array the turning point $x^*$ places the absorption of each frequency, and the damping sets the width of the absorbing zone. Grading the resonator mass with $\omega_r$, equation~\eqref{eq:mass_grading_rule}, makes this width nearly uniform: with exponential grading and $\sqrt{\mu^*}=0.163$ it varies by $13\%$ over $\omega\in[3,7]~\si{\radian\per\second}$, against a factor of $3.2$ with uniform loading (Section~\ref{subsec:num_graded}). A density taper, equation~\eqref{eq:density_grade}, reflects less than an abrupt array for $L_g\ge5\lambda$, but only modestly: the band-averaged $|R|^2$ is $8$--$23\%$ lower without damping and $1$--$21\%$ lower at $\zeta=0.02$, and the gain shrinks as $L_g$ grows. The length of the array matters more than its taper: at $\zeta=0.02$ the absorbed fraction rises from $0.26$ for $L_g=\lambda$ to $0.88$ for $L_g\ge20\lambda$, for both arrays.

\subsection{Numerical aspects}\label{subsec:disc_numerics}

\paragraph{Truncation.} Sponges that act on the structure do not damp the evanescent fluid modes, which carry no energy flux. Sources and scatterers are therefore kept at least $5/\operatorname{Im}k_1$ from the sponges, about $32~\si{\meter}$ at $\omega=10~\si{\radian\per\second}$ and $H=10~\si{\meter}$, so that the reflected evanescent field is reduced by $\mathrm{e}^{-10}$ (Section~\ref{subsubsec:num_setup_A}). In configuration~B the open water on each side is at least $6/\kappa_1$ long (Section~\ref{subsubsec:num_setup_B}). Bending sponges are also transparent to long waves; at $\omega\le1~\si{\radian\per\second}$ the waves are fully reflected. Closed channels then have basin resonances at low frequency. The channel length is chosen midway between them, $L=(n+\tfrac14)\lambda_0$, and at $\omega_r=1~\si{\radian\per\second}$ the bending sponges are replaced by viscous dashpot layers, which absorb all wavelengths (Section~\ref{subsubsec:num_L1_3b}). In addition, all coefficients are extracted with relations that are exact for any amount of reflection from the domain ends.

\paragraph{Attachments and polynomial degree.} A resonator inside an element lowers the accuracy of $p=3$. Its resonance frequency is then reproduced to $0.7$--$1.8\times10^{-6}$ at order $\approx3$, against $3\times10^{-11}$--$3\times10^{-9}$ on a node (Section~\ref{subsubsec:num_L1_3}). Resonators are therefore placed on mesh nodes. On the three-dimensional plate the location within a cell changes the driving-point compliance by at most $0.20\%$ on the finest mesh, and the difference falls fourfold per refinement (Section~\ref{subsubsec:num_L3_2}). Band edges at small spacing need $p=3$ and meshes sized with the cell rather than with the wavelength. At $a=0.25~\si{\meter}$ the gap is $7.7\times10^{-5}\,\omega_r$, so the edges must be accurate to $10^{-7}$ (Section~\ref{subsubsec:num_L1_6}). With $p=3$ the gap is reproduced to between $10^{-11}$ and $6\times10^{-9}$ at every spacing; with $p=2$ and $64$ elements per cell, to $1.5$--$3\times10^{-4}$. For low-frequency resonators the accuracy of the renormalised resonance is limited by the resolution of the flexural near field, not of the wavelength, so meshes should be refined around the attachments.

\paragraph{Sources for finite plates.} Plane waves incident from open water on a free plate are almost totally reflected at the free edge: $|R|>0.99$ at the frequencies of Section~\ref{subsec:num_graded}, and almost no power enters the plate. This limits the power that open-water waves can deliver to resonators on a long free plate. For the graded arrays the incident plate wave is therefore generated by a point source on a bare lead-in, with viscous sponges of reflection at most $1.4\times10^{-3}$ at both ends of the plate. The net plate-wave flux into the array then balances the dashpot power exactly, to $9.3\times10^{-7}$ in the computations.

\paragraph{Reuse of factorisations.} Two devices avoid repeated solves. Bloch-phased loading of periodic supercells gives band frequencies from the forced-response solver, without quasi-periodic boundary conditions or an eigenvalue solver (Section~\ref{subsubsec:num_setup_lattice}). The Woodbury identity~\eqref{eq:lat_woodbury} holds to round-off, so any resonator set is evaluated from one bare factorisation and an $N\times N$ solve per frequency (Section~\ref{subsubsec:num_L2_1}).

\paragraph{Three-dimensional cost.} In three dimensions the sparse direct factorisation dominates both memory and time. A sizing study on slab-shaped meshes like those of Section~\ref{subsec:num_3D} [L3-0] shows that the factor fill grows as $N^{1.25}$--$N^{1.5}$ with the number of unknowns $N$, depending on the ordering. A nested-dissection (METIS) ordering of the LU factorisation reduces the fill by a factor of $2.4$ and the operation count by $7.6$ at $4.7\times10^5$ unknowns, and both gains grow with $N$ [L3-0]. With this ordering the runs of Sections~\ref{subsubsec:num_L3_1} and~\ref{subsubsec:num_L3_2}, up to $3.8\times10^5$ unknowns, need at most $12$~GB of memory. Even so, a single resonator in a quarter domain or an anisotropic island in a half domain with buffers of $5/\kappa_1$ exceeds $20$~GB [L3-0], so those configurations need reduced buffers, a cluster or an iterative solver.

\paragraph{Driving-point values in a channel.} The three-dimensional point resonator needs no footprint: the local field follows the biharmonic Green function, with the $r^2\ln r$ coefficient within $2.3\%$ (resonator force) and $3.7\%$ (point load) of its exact value (Section~\ref{subsubsec:num_L3_2}). In a strip narrower than $\lambda_p$, $G_h$ is a global quantity of the strip, so the far-field cells must stay fixed while the attachment is refined. With a fixed far field the driving-point compliance converges at rates of $1.6$--$1.9$ and is resolved to about $0.3\%$, against about $5\%$ when the far-field cells change with the refinement. A residual increment of $0.3\%$ on the finest mesh is attributed to elongated cells of the tensor-product grading.

\subsection{Limitations}\label{subsec:disc_limitations}

The analysis is linear: potential flow, small deflections and linear mass--spring--dashpot resonators, whose dashpot is an idealised linear power take-off. Nonlinear wave loading, large structural motion and nonlinear power take-off are outside its scope. The results are computed in the frequency domain, although the formulation of Section~\ref{sec:numerical_formulation} also holds in the time domain.

The resonators are point attachments, i.e.\ line attachments per unit width in the vertical section. The local theory replaces them by a continuous field, and the discrete model by point forces. Section~\ref{subsubsec:num_L3_2} shows that a point attachment on a Kirchhoff plate needs no footprint, but the finite size and the dynamics of real attachments are not modelled.

Most results are for one plate and one resonator family (Table~\ref{tab:dispersion_test_parameters}) and for line arrays in the vertical section. The homogenised resonator field is tested in the vertical section only (Section~\ref{subsubsec:num_L1_7}). The three-dimensional results are limited by the cost of direct solvers. The bare plate under oblique incidence reproduces the isotropic dispersion relation and the oblique Green function, with errors below those at normal incidence (Section~\ref{subsubsec:num_L3_1}). The single resonator and its absorption cross-section (L3-3), the anisotropic resonator island of Section~\ref{subsubsec:aniso_rays} (L4-4) and the forced response of a finite plate against equation~\eqref{eq:mode_dilution} (L4-8) are being completed [PENDING]. Until then, the anisotropic predictions of Section~\ref{subsec:anisotropic_density} and the three-dimensional mode dilution remain estimates.

\section{Conclusions}\label{sec:conclusions}

\paragraph{Formulation.} We have formulated a monolithic finite element model for linear potential flow coupled to a floating pre-tensioned plate, discretised with a C/DG method, and to discrete mass--spring--damper resonators attached at points. The resonator amplitudes are global scalar unknowns, and the attachments are evaluated by Dirac measures at any point of the structure. The discrete problem satisfies an energy identity and a power balance, and it is reciprocal: reciprocity holds to $1.7\times10^{-9}$ or better in the computations. Eliminating the resonators is exact at the discrete level, so the transfer compliance $G_B$ of the bare structure and the Woodbury identity evaluate any resonator set from one factorisation. The wet modal analysis converges in three fixed-point iterations and agrees with the resonance written with $G_B$ to $0.7$--$1.3$ half-widths of the resonance.

\paragraph{Local theory.} The local effective-medium dispersion relation of a floating plate with resonators shows that the fluid dilutes the resonators. The mass ratio $\mu(k)$ depends on the wavenumber and vanishes at long wavelength. The hybridisation at the crossing is $\sqrt{\mu(k_c)}$, and the resonator mass must be compared with a fluid layer of thickness $1/k_c$. The homogenised floating plate has no complete stop band. The optimal damping of a single resonator collapses at low frequency, by more than two decades between $\omega_r=10$ and $1~\si{\radian\per\second}$. For graded arrays, the turning point locates the absorption, its width is set by the damping, mass grading makes it nearly uniform over most of the band, and density tapers reduce the reflection only modestly.

\paragraph{Discrete lattices.} Lattice sums and multiple scattering show that the complete stop band of a floating lattice is a discreteness effect. Its upper edge is the resonance renormalised by the lattice near field, its lower edge comes from the zone-folded component, and at fixed loading its width vanishes as $a^4$. Long floating arrays have a weakly attenuated Bloch mode that continues the slowest evanescent fluid mode and bounds their attenuation per unit length by about $\pi/(2H)$. On a finite structure, a single resonator can absorb more than one half of the incident power where the structure radiates asymmetrically, up to $91\%$ in our examples. The hybridisation available on a finite structure decreases with its size.

\paragraph{Finite element evidence.} The finite element model shares no part of its solution with the Green-function, lattice-sum and Foldy references, and reproduces them to the accuracy of the discretisation. Smooth quantities converge at order $2(p-1)$. With $p=3$ the band edges and the hybridisation of the line lattice agree with the lattice sums to $5\times10^{-9}$ or better, and the finite-array coefficients agree with the Foldy solution to below $10^{-6}$ on the finest meshes. In configuration~B the Woodbury identity holds to round-off, and the absorption bound is reached to $1.8\times10^{-4}$. The computations also correct the local estimates where they fail: the width of graded absorption, the prefactor of the size law and the range of validity of the near-field estimate. In three dimensions, the bare plate reproduces the isotropic dispersion relation under oblique incidence, and a point resonator on the plate needs no footprint [PENDING: summary of L3-3, L4-4 and L4-8].
}
